\section{The minimal marked continuation quotient}
\label{sec:v10-marked}

We now address exact state reduction. A statistic sufficient for the
parameter need not be sufficient for a latent decision target. The
correct terminal datum is the joint conditional law of the parameter
and the designated mark. We give both an intrinsic operational
characterization and a finite construction.

Consider a finite hidden instrument with initial law $\nu_\theta$,
transition/report probabilities $K_{\theta,t}(z',y\mid z,a)$, and a
terminal mark channel $L_\theta(w\mid z_T)$ on a finite set $\mathsf W$.
Here $K$ uses the new state $z'$ and the report jointly; no old-state
observation channel is silently substituted. A finite collection of
path targets can be included by adjoining their necessary finite
memory to $z$. This augmentation, when used, is part of the input size.
Fix a reference prior $\pi(\theta)>0$. For every supported controlled
history $h=(a_0,y_1,\ldots,a_{t-1},y_t)$, let
$\alpha_t(h)$ be the posterior row on $(\Theta,Z_t)$ under the
reference prior and the indicated past actions. Control probabilities
cancel in Bayes' formula; histories of zero reference probability are
assigned an arbitrary existing label. No extra retained label is needed
for events that have probability zero under every parameter.

\begin{definition}[Continuation experiment]\label{def:v10-continuation}
For each future open-loop action word $v=a_t\cdots a_{T-1}$, let
$\mathcal C_t(h;v)$ be the conditional law of
$(\Theta,Y_{t+1:T},W)$ given $h$, with future actions $v$.
Two supported histories are equivalent, written $h\equiv_t h'$, if
$\mathcal C_t(h;v)=\mathcal C_t(h';v)$ for every such $v$.
Denote their class by $c_t(h)$ and the number of classes by $K_t$.
\end{definition}

This definition compares experiments rather than lists of selected
coordinates. At $T$ it reduces to equality of
$\mathcal L(\Theta,W\mid h)$, not only equality of
$\mathcal L(\Theta\mid h)$. Since the parameter set is finite and the
prior is strictly positive, changing the reference prior applies an
invertible positive reweighting to every continuation law. The
classes are consequently independent of this choice.

\begin{lemma}[Controlled continuation]\label{lem:v10-continuation}
The classes satisfy
\[
 h\equiv_t h'\ \Longrightarrow\
 hay\equiv_{t+1}h'ay
\]
for every $a,y$ of positive conditional probability. The probability
of $y$ is itself constant on a class. Equality of the open-loop
continuation laws implies equality under every randomized feedback
policy using the subsequent report history.
\end{lemma}
\begin{proof}
The marginal of the first future report is a marginal of any
continuation law beginning with $a$. It is the same for $h,h'$.
Condition the joint law of $(\Theta,Y_{t+1:T},W)$ on that report
being $y$. Its positive common denominator is the first-report
probability, and its remaining numerator is the continuation law
from $hay$ or $h'ay$. Varying the remaining action word gives the
successor assertion. For a feedback policy, multiply each open-loop
path probability by the product of its conditional control
probabilities along that path and sum over action words. These
factors depend only on the displayed future reports and on independent
randomization. Equality is preserved. Null successors have zero
probability under every parameter of positive reference prior.
\end{proof}

\subsection{An operational necessity theorem}
The category in which a minimality claim is made matters. At a
checkpoint $t$, a \emph{restart preparation} chooses a supported past
history $H$ with an arbitrary law $\eta$, and then runs its conditional
continuation experiment. The candidate statistic is obtained by
consuming that same history through a channel $Q_t(dm\mid h)$.
It does not resample a new history or replace the actual target $W$.
A restart prediction task fixes a future action word $v$ and an event
$A\subseteq\Theta\times Y_{t+1:T}\times\mathsf W$. Immediately at the
checkpoint, before the continuation is observed, the decision maker
predicts the indicator of $A$ with an action in $[0,1]$ and squared
loss. The full-history optimal prediction is
$p_A^v(h)=\mathcal C_t(h;v)(A)$.

Call a source-consuming causal statistic \emph{restart-decision
complete} when it preserves the full-history Bayes risk in every one
of these tasks, for every preparation $\eta$ and every checkpoint.
The definition requires equality of operational risks. It does not
require any posterior coordinate to be explicitly stored.

\begin{theorem}[Marked operational minimality]\label{thm:v10-minimal}
A privately randomized finite-valued statistic $M_t$ is restart-decision complete at $t$
if and only if $c_t(H)$ is recoverable as a function of $M_t$ on every
supported history. In particular any complete randomized causal
machine has at least $K_t$ retained states at checkpoint $t$.
The deterministic quotient $(c_t)$ is simultaneously complete at all
checkpoints, has exactly these widths, and is minimal in this category.
For a public seed $U$, the corresponding criterion is recovery of
$c_t(H)$ from $(U,M_t)$, not from $M_t$ alone. The same lower bound
on the number of private register labels holds in either signature.
\end{theorem}
\begin{proof}
The optimal Brier risks before and after compression differ by
\begin{equation}\label{eq:v10-brier}
 \E_\eta\left[
   \bigl(p_A^v(H)-\E_\eta[p_A^v(H)\mid M_t]\bigr)^2\right].
\end{equation}
Indeed $\E[\1_A\mid H,M_t]=p_A^v(H)$, since the statistic's
randomization is independent of the actual continuation given $H$.
Orthogonal projection in $L^2$ proves the identity. There are only
finitely many histories and cylinder events in the present finite
experiment. Choose a preparation with positive mass on every
supported history. Vanishing of \eqref{eq:v10-brier} for every $v,A$
forces $p_A^v(h)=p_A^v(h')$ whenever both histories give positive
probability to a common value $m$. Thus $h\equiv_t h'$. Each retained
value has support in one equivalence class, so the class is a function
of that value and there must be at least $K_t$ values.

Conversely, if the class is recoverable, every $p_A^v(H)$ is measurable
from $M_t$, and \eqref{eq:v10-brier} is zero for every preparation.
Lemma~\ref{lem:v10-continuation} implements the classes by deterministic
causal updates. It also proves equality for feedback continuations.
For an independent public seed $U$, apply the same argument to
$(U,M_t)$; if the declared width bounds only $M_t$, a free retained
$U$ is a genuinely different signature. In an exact complete public
scheme with finite widths, zero nonnegative excess also holds for
almost every fixed seed, so the lower bound $S_t\ge K_t$ remains
valid. One cannot use a seed to merge distinguishable histories while
preserving all these zero-loss differences.
\end{proof}

\begin{proposition}[A quantitative separation]\label{prop:v10-overlap}
Suppose a continuation event has probabilities differing by $\Delta$
at $h,h'$. For the equal two-history preparation, any randomized
statistic has Brier excess at least
\[
 \frac{\Delta^2}{4}
 \bigl(1-\TV(Q_t(\cdot\mid h),Q_t(\cdot\mid h'))\bigr).
\]
\end{proposition}
\begin{proof}
Write $x_m=Q_t(m\mid h)$ and $y_m=Q_t(m\mid h')$. Direct evaluation
of the conditional variance in \eqref{eq:v10-brier} gives
$\frac{\Delta^2}{2}\sum_m x_my_m/(x_m+y_m)$, with zero terms omitted.
Since $xy/(x+y)\ge\min(x,y)/2$ and
$\sum_m\min(x_m,y_m)=1-\TV(x,y)$, the result follows.
\end{proof}

This is not a minimality assertion for an arbitrary independent
simulator with unpriced memory. Nor is it an assertion that a chosen
small list of losses separates every latent feature. The restart
family is decision complete for the stated marked continuation
experiment; its loss tests force the necessary state rather than
assuming a posterior coordinate representation in the definition.
The continuation and linear-representation aspects have classical
antecedents in automata and predictive-state theory
\cite{Nerode1958,ShaliziCrutchfield2001,Kiefer2020}; the target mark and
the source-consuming operational category are essential here.

\subsection{Two exact marked causal channels}
Let $k^0_t(y\mid c,a)$ be the reference next-report probability and
$F_t(c,a,y)$ the successor class. Define
\[
 H^0_t(c'\mid c,a)=
 \sum_{y:F_t(c,a,y)=c'}k^0_t(y\mid c,a),\qquad
 \lambda_\theta(c)=\mathbb P_\pi(\Theta=\theta\mid c).
\]
On the support of parameter $\theta$, Bayes' formula gives
\begin{equation}\label{eq:v10-ratio}
 k_{\theta,t}(y\mid c,a)
 =k^0_t(y\mid c,a)
   \frac{\lambda_\theta(F_t(c,a,y))}{\lambda_\theta(c)}.
\end{equation}
The marked terminal channel is
$G_\theta(w\mid c_T)=\mathbb P_\pi(W=w\mid\Theta=\theta,c_T)$.
It is part of the quotient; it is not discarded.

\begin{theorem}[Exact marked reconstruction]\label{thm:v10-reconstruction}
The original instrument and the class instrument admit causal
channels in both directions preserving the joint law of the reported
transcript and the actual terminal mark, under every parameter and
admissible feedback policy. The reverse report kernel is
\begin{equation}\label{eq:v10-reconstruct}
 L_t(y\mid c,a,c')=
 \frac{\1_{\{F_t(c,a,y)=c'\}}k^0_t(y\mid c,a)}
      {H^0_t(c'\mid c,a)}
\end{equation}
when the denominator is positive. It is independent of $\theta$.
The class instrument uses transitions
$H_{\theta,t}(c'\mid c,a)
=H^0_t(c'\mid c,a)\lambda_\theta(c')/\lambda_\theta(c)$
and terminal channel $G_\theta$. Thus the channels preserve every
bounded loss on the declared marked experiment, not just parameter
likelihoods.
\end{theorem}
\begin{proof}
The forward channel consumes each actual report and updates $c$ by
$F_t$; its current class is a finite register. For the reverse channel,
first observe the new class $c'$, then draw $y$ by
\eqref{eq:v10-reconstruct}. Multiplying this kernel by
$H_{\theta,t}$ gives exactly \eqref{eq:v10-ratio}. The terminal
conditional distribution of $W$ given the full original history and
$\theta$ is $G_\theta(\cdot\mid c_T)$ by the definition of the terminal
class. Consequently the product of transition factors, report
reconstruction factors and this terminal mark factor is the original
joint density of $(Y_{1:T},W)$, for each open-loop action word.
The feedback factors can be inserted before summing action words,
as in Lemma~\ref{lem:v10-continuation}. This proves the claimed joint,
parameter-uniform identity. Cases with zero $\lambda_\theta$ or zero
transition probability lie outside that parameter's support and may
be assigned arbitrary kernels. The positive reference prior makes
these conventions consistent across parameters.

A reverse transducer needs the previous class, the new class and the
current action while forming its report. Between updates it retains
only the previous class. Composing it with an $S_t$-state machine is
therefore possible with at most $K_tS_t$ retained labels at time $t$;
the new input is transient. These resource costs are separate from
the bare exact-simulation identity.
\end{proof}

\begin{example}[A parameter-sufficient statistic that loses its mark]
\label{ex:v10-xor}
Let $\mathbb P(\Theta=0)=9/10$, let $Y$ be an independent fair bit,
and put $W=Y\mathbin{\oplus}\Theta$. The parameter posterior is
constant in $Y$. A terminal quotient that keeps only that posterior
has one class. Yet the Bayes error in guessing the actual $W$ is
$1/10$ with $Y$ and $1/2$ with that class alone. The marked terminal
laws $\mathcal L(\Theta,W\mid Y=0)$ and
$\mathcal L(\Theta,W\mid Y=1)$ differ, so the marked quotient has two
classes. Its reverse channel preserves the relation
$W=Y\mathbin{\oplus}\Theta$, not merely the marginal fair-bit law of
$Y$. This example explains why an exact reconstruction theorem for
reports alone cannot be used for arbitrary latent-target losses.
\end{example}
